\begin{frontmatter}

\title{Member-specific thresholds and contested-price intervals for stage-gate screening: A curvature-based real-options method under regulatory ambiguity}

\begin{abstract}
Stage-gate committees decide when to make irreversible investments under regulatory ambiguity, and members weighting adverse scenarios differently can reach opposite decisions at the same price. Real-options analysis reduces this decision to one project-level trigger, and committee decision support compresses member assessments into a score or ranking, so neither reports the price range over which member classifications differ. This study therefore derives one price threshold per member from affine combinations of the triggers at both ends of a volatility interval, so a committee can locate disagreement before any vote. The extreme thresholds bound a contested-price interval, and a unique critical weight splits members around the no-ambiguity anchor. Under the affine rule, trigger curvature places this weight above one half in standard-deviation units and below it in variance units. Against a pre-commitment $\alpha$-maxmin expected utility benchmark, the largest relative trigger error over the primary radii was 2.74\%. Because small trigger errors did not guarantee identical classifications, affine-only release is capped at 10\%, the largest evaluated radius at which two fidelity criteria held. In Monte Carlo committee panels, the member-threshold vector retained investment-timing information, with 88.5\% lower median loss than the ex post best of four single-trigger summaries. In a carbon capture, utilization, and storage retrofit, all archived prices stayed below the lower edge under the reference configuration. In 15 of 32 sensitivity configurations, however, the highest archived price reached or exceeded this edge. Unanimity at a gate therefore depended on the recorded inputs, and the member-level report made this dependence visible.
\end{abstract}

\begin{keyword}
Real options \sep Stage-gate screening \sep Regulatory ambiguity \sep Group decision support \sep CCUS investment
\end{keyword}

\end{frontmatter}

\section{Introduction}

Large, irreversible investments are commonly reviewed through a sequence of stage gates \citep{Cooper2008}. At each gate, a committee decides whether to commit to the project now or to defer it to a later review. Real-options analysis reduces this timing decision to a single investment trigger \citep{McDonaldSiegel1986,DixitPindyck1994}. Commitment is warranted once project value reaches this trigger. A single trigger gives the review one compact criterion, which can be documented and recomputed at successive gates. When project cash flows are linked to an observable price, the trigger can be restated on the price axis, in units the committee already monitors. In this paper, a trigger is a project-value quantity, and a threshold is its image on the price axis.

Committee members, however, may receive the same engineering, cost, and policy evidence and still place different weights on adverse regulatory scenarios. An ambiguity-aware real-options analysis returns one project-level trigger under a representative ambiguity attitude, so member-level dispersion disappears at the modeling stage. Committee decision-support procedures can retain individual assessments, but they aggregate them into a score, a ranking, or a vote before reporting. In both cases, the assessments are compressed before their implied classifications are compared.

This compression matters for the gate only when member-level dispersion changes the classification, which does not always happen. All members receive a wait classification at a price below every member threshold and an invest classification at a price at or above every threshold. A mixed classification arises only between the lowest and the highest member thresholds, which bound the contested-price interval. The position of the price relative to the whole threshold set therefore determines whether one classification suffices. An aggregate score replaces this set with one number and thus hides whether the price lies between the extreme thresholds. The set, however, remains recoverable from the recorded calibration, ambiguity radius, and member weights, and changing these inputs can reclassify the same prevailing price.

This paper determines when member heterogeneity changes the current gate classification and reports the answer without aggregating member thresholds in advance. All members share one ordered ambiguity set and differ only in their recorded weights on its adverse endpoint. Member triggers follow a closed-form affine endpoint rule, and the $\alpha$-maxmin expected utility ($\aMEU$) optimum under the same inputs serves as its accuracy benchmark. The two rules can classify the same price differently. The paper therefore also determines when the affine report can be released alone and when the benchmark must be recomputed.

The shared ambiguity set represents one component of regulatory uncertainty as an ordered interval of forward volatilities, with a declared radius around an anchor. Its adverse endpoint is the low-volatility endpoint, because lower volatility reduces the value of the option to wait and lowers the trigger.

The affine rule also determines whether each member trigger lies above or below the no-ambiguity anchor. The trigger is strictly increasing in volatility, so at a positive radius the anchor trigger lies strictly between the two endpoint triggers. Exactly one weight in $[0,1]$, called the critical weight, therefore places a member trigger at the anchor. This weight splits the members into two groups around the anchor, and the split carries over to the price thresholds. The critical weight is located relative to one half, the weight at which the two endpoint triggers count equally.

The running example is a committee screening a carbon capture, utilization, and storage (CCUS) retrofit of a coal-fired power plant. The retrofit costs $1.585\times10^{9}$\,CNY and captures $1.69\times10^{6}$ tonnes of CO$_2$ per year \citep{Sun2024}. Project cash flows track an observable carbon price, so the member triggers become marginal carbon-price thresholds in Chinese Yuan per tonne (CNY/t). On the input side, the ambiguity radius of each gate comes from a policy-text series.

The paper makes one main contribution and two supporting contributions. The main contribution is a member-indexed gate report, which consists of the member-threshold vector and the contested-price interval bounded by its extremes. The report provides a pre-vote test of whether the prevailing price receives the same classification under every member assessment. The two supporting contributions are analytical and methodological. The analytical contribution is a location rule for the critical weight. Strict monotonicity of the affine rule makes this weight unique, and trigger curvature places it relative to one half. This curvature changes sign between standard-deviation and variance coordinates, so the direction of shift for members weighted near one half depends on the declared coordinate. The methodological contribution is a release rule for the affine report. Benchmark comparisons and a classification-fidelity analysis measure the accuracy of this report in price units. In this analysis, small trigger errors could still change the invest-or-wait outcome. A finite-registry safeguard converts these accuracy limits into conditions for affine-only release, and each release can then be audited against the registered configurations.

The CCUS application produced two main results. First, unanimity depended on the configuration. Under the reference configuration, every archived price stayed below the contemporaneous lower interval edge. By contrast, in 15 of the 32 sensitivity configurations, the highest archived price reached or exceeded this edge. Second, the common appraisal inputs moved the interval edges further than member disagreement did. A 10\% reduction in the cost-to-volume ratio $I/Q$ moved the lower edge by 11.34\,CNY/t. By comparison, the widest interval under the reference profile along the reconstructed path was 6.63\,CNY/t.

By treating member disagreement as a reported quantity with stated accuracy limits, the study addresses a reporting gap between real-options analysis and committee decision support. This gap separates a trigger computed for the project from a decision taken by its committee.
